% ==STYLE-BEGIN==
\documentclass[UTF8,11pt]{ctexart}
\usepackage{amsmath,amssymb,amsthm}
\usepackage[a4paper,margin=2.4cm,top=2.3cm,bottom=2.3cm]{geometry}
\usepackage{fancyhdr}
\usepackage{enumitem}
\usepackage{hyperref}
\hypersetup{hidelinks}

% ---- 版式 ----
\linespread{1.15}
\setlength{\parskip}{0.35em}
\setlength{\parindent}{1.6em}
\setlist{nosep,leftmargin=2em}

% ---- 页眉 ----
\pagestyle{fancy}
\fancyhf{}
\fancyhead[L]{\small\coursename}
\fancyhead[R]{\small\thepage}
\renewcommand{\headrulewidth}{0.4pt}

% ---- 定理环境：中文用正体，标题加粗，按小节共用一条流水号 ----
\newtheoremstyle{cnthm}{0.7em}{0.7em}{\normalfont}{}{\bfseries}{}{0.5em}{}
\theoremstyle{cnthm}
\newtheorem{thm}{定理}[section]
\newtheorem{lem}[thm]{引理}
\newtheorem{prop}[thm]{命题}
\newtheorem{cor}[thm]{推论}
\newtheorem{defn}[thm]{定义}
\newtheorem{rem}[thm]{注}
\newtheorem{exa}[thm]{例}
\newtheorem{conj}[thm]{猜想}
\renewcommand{\proofname}{\bfseries 证明}
\providecommand{\coursename}{}
% ==STYLE-END==
\renewcommand{\coursename}{数学分析\,·\,1102}
\title{数学分析 1102}
\date{}
\begin{document}
\maketitle

\section{微分(续)}

\begin{rem}[Recall: 可微]
设 $f(x)$ 在 $x_0$ 的一个开区间内有定义. 若在此开区间内
\[
f(x)=f(x_0)+A(x-x_0)+o(x-x_0),
\]
则称 $f$ 在 $x_0$ 处\textbf{可微}.
\end{rem}

\begin{thm}
$f(x)$ 在 $x_0$ 处可微 $\iff$ $f(x)$ 在 $x_0$ 处可导, 且 $A=f'(x_0)$.
\end{thm}

\begin{defn}[微分]
称
\[
\mathrm{d}f:=f'(x_0)\,\Delta x
\]
为 $f$ 在 $x_0$ 处的\textbf{微分}.
\end{defn}

\begin{rem}
(1) 当 $x_0$ 固定时, $\mathrm{d}f$ 是 $\Delta x$ 的函数;

(2) $f'(x_0)=\dfrac{\mathrm{d}f}{\mathrm{d}x}$ 是一个"商"(微分之商);

(3) 过点 $\bigl(x_0,f(x_0)\bigr)$、以 $f'(x_0)$ 为斜率的直线就是 $f$ 在 $\bigl(x_0,f(x_0)\bigr)$ 处的切线:
\[
y-f(x_0)=f'(x_0)(x-x_0)=f'(x_0)\,\Delta x=\mathrm{d}f(\Delta x);
\]

(4) 实质: 用线性的切线近似 $f(x)$, 其差值是高阶无穷小;

(5) $f'(x_0)$ 就是 $f$ 的微分 $\mathrm{d}f$ 与自变量的微分 $\mathrm{d}x\triangleq\Delta x$ 之商.
\end{rem}

\section{微分形式的不变性}

\begin{thm}[微分形式的不变性]
设 $y=y(u)$ 是关于 $u$ 的函数, $u=u(x)$ 是关于 $x$ 的函数. 则 $\mathrm{d}y$ 既可以看成 $u$ 的微分, 也可以看成 $x$ 的微分, 而两种算法结果相同:
\[
\mathrm{d}y=y'(u)\,\mathrm{d}u
\qquad\text{与}\qquad
\mathrm{d}y=y'(x)\,\mathrm{d}x .
\]
\end{thm}

\begin{proof}
因为 $u=u(x)$ 可微, 其微分为
\[
\mathrm{d}u=u'(x)\,\mathrm{d}x .
\]
代回第一个式子:
\[
\mathrm{d}y=y'(u)\,\mathrm{d}u=y'(u)\,u'(x)\,\mathrm{d}x=y'(x)\,\mathrm{d}x ,
\]
与第二个式子一致.
\end{proof}

\begin{rem}[两处 $\mathrm{d}u$ 不是一回事]
上面两个式子里出现的 $\mathrm{d}u$ 含义不同:

(1) 在 $\mathrm{d}y=y'(u)\,\mathrm{d}u$ 中, 若 $u$ 是\textbf{自变量}, 则按定义 $\mathrm{d}u$ 就取增量 $\Delta u$ (自变量的微分等于它的增量);

(2) 在 $u=u(x)$ 的情形, $\mathrm{d}u=u'(x)\,\mathrm{d}x$ 是 $u$ 作为 $x$ 的函数的\textbf{微分} (线性主部), 它与增量
\[
\Delta u=u(x+\mathrm{d}x)-u(x)=u'(x)\,\mathrm{d}x+o(\mathrm{d}x)
\]
相差一个 $o(\mathrm{d}x)$.

所以两个 $\mathrm{d}u$ 不是同一个量: 前者是 $\Delta u$, 后者是真微分, 二者差 $o(\mathrm{d}x)$. 微分形式不变性说的正是: 无论 $u$ 是自变量还是中间变量, $\mathrm{d}y=y'(u)\,\mathrm{d}u$ 这个\textbf{形式}都成立, 只是 $\mathrm{d}u$ 的含义随情形而变; 这也正对应复合函数求导的链式法则
\[
\frac{\mathrm{d}y}{\mathrm{d}x}=\frac{\mathrm{d}y}{\mathrm{d}u}\cdot\frac{\mathrm{d}u}{\mathrm{d}x} .
\]
\end{rem}

\section{近似计算}

由 $f(x)=f(x_0)+f'(x_0)(x-x_0)+o(x-x_0)$ 知, 当 $x$ 与 $x_0$ 较接近时
\[
f(x)\approx f(x_0)+f'(x_0)(x-x_0).
\]

\begin{exa}
求 $\sqrt[n]{1+x}$ 在 $|x|$ 足够小时的近似式.

令 $f(x)=\sqrt[n]{1+x}$, 取 $x_0=0$, 则 $f(x)-f(0)\approx f'(0)(x-0)$. 而
\[
f'(x)=\frac1n(1+x)^{\frac1n-1}\quad\Longrightarrow\quad f'(0)=\frac1n ,
\]
故
\[
f(x)\approx f(0)+\frac1n x=1+\frac{x}{n},
\qquad\text{即}\qquad
\sqrt[n]{1+x}\approx 1+\frac{x}{n}.
\]
\end{exa}

\begin{exa}
估算 $\sqrt{420}$.
\[
\sqrt{420}=20\sqrt{\frac{420}{400}}=20\sqrt{1+\frac1{20}}
\approx 20\left(1+\frac1{2}\cdot\frac1{20}\right)=20\times1.025=20.5 .
\]
\end{exa}

\begin{exa}
设 $y=\ln(1+e^{x})$, 求 $\mathrm{d}y$.

令 $y=\ln u$, $u=1+e^{x}$, 则
\[
\mathrm{d}y=y'(u)\,\mathrm{d}u=\frac{\mathrm{d}u}{u}
=\frac{e^{x}\,\mathrm{d}x}{1+e^{x}}
=\frac{e^{x}}{1+e^{x}}\,\mathrm{d}x .
\]
\end{exa}

\end{document}
